\chapter{Programación Funcional}
\label{ch:programacion_funcional}

La programación funcional representa una ruptura fundamental con el modelo de computación de Von Neumann. Mientras que el paradigma imperativo se centra en \textbf{cómo} modificar el estado de una máquina mediante asignaciones sucesivas, el paradigma funcional se enfoca en \textbf{qué} es el valor computado mediante la aplicación de funciones puras \cite{SICP}. Este enfoque tiene sus raíces en el Cálculo Lambda ($\lambda$-calculus) desarrollado por Alonzo Church en la década de 1930, el cual demostró ser un modelo de computación universal equivalente a la Máquina de Turing \cite{PFPL, TAPL}.

En este capítulo, exploraremos cómo la ausencia de estado mutable y el tratamiento de las funciones como «ciudadanos de primera clase» permiten construir software más robusto, compacto y fácilmente paralelizable, utilizando el lenguaje \eqk{Scheme} como nuestro vehículo de experimentación \cite{HtDP, SICP}.

\section{Filosofía y Fundamentos del Paradigma Funcional}
\label{sec:4_1_filosofia}

El corazón de la programación funcional es la noción de que un programa es una descripción de una correspondencia matemática entre un dominio de entrada $\eqdom{D}$ y un rango de salida $\eqdom{R}$ \cite{PFPL}.

\subsection{La Naturaleza de la Programación Funcional}
\label{subsec:4_1_1_naturaleza}

A diferencia de los lenguajes tradicionales donde el programador gestiona explícitamente la memoria y el flujo de control, un lenguaje funcional puro opera bajo un conjunto de restricciones que, paradójicamente, aumentan su poder expresivo \cite{HtDP, SICP}.

\begin{definicion}{Programación Funcional}
Un lenguaje de programación funcional se define por las siguientes propiedades core \cite{HtDP}:
\begin{itemize}
    \item \textbf{Ausencia de Efectos Secundarios:} Las funciones no modifican variables externas ni producen cambios observables en el mundo exterior (como operaciones de I/O impredecibles) más allá de su valor de retorno \cite{SICP, HtDP}.
    \item \textbf{Inmutabilidad (Sin Asignaciones):} No existe el operador de asignación tradicional ($=$). Una vez que un nombre se liga a un valor $\equ{x} \eqop{:=} \eqk{v}$, ese valor permanece constante durante todo su ámbito de validez \cite{SICP}.
    \item \textbf{Transparencia Referencial:} El resultado de evaluar una función $\eqfn{f}(\equ{x})$ depende únicamente de sus argumentos. Esto implica que cualquier llamada a una función puede ser sustituida por su valor resultante sin alterar el comportamiento del programa \cite{SICP}.
\end{itemize}
\end{definicion}

Desde una perspectiva de ingeniería, este paradigma permite que el programador se dedique casi exclusivamente al diseño del algoritmo. Se estima que en lenguajes como \eqk{C} o \eqk{Java}, hasta el 50\% del código está dedicado a la gestión de estados y memoria, mientras que en los lenguajes funcionales, el 100\% del código suele representar la lógica del dominio \cite{HtDP}.

\begin{observacion}{Código más compacto y correcto}
La programación funcional facilita la verificación formal y el razonamiento matemático sobre el código. Al eliminar las variables globales y los estados ocultos, se reduce drásticamente la probabilidad de introducir errores de tipo «condición de carrera» o efectos colaterales imprevistos, lo que resulta en un código significativamente más conciso \cite{SICP, HtDP}.
\end{observacion}

\begin{fisicaconexion}{Sistemas conservativos y reversibilidad}
En física, un sistema conservativo es aquel donde la energía total permanece constante y el estado final depende únicamente de las condiciones iniciales. La programación funcional emula esta naturaleza: al ser referencialmente transparente, el «estado» informativo se conserva y transforma de manera predecible. La ausencia de efectos secundarios es análoga a un sistema sin fricción, donde no hay pérdida de energía (información) hacia entornos no modelados, permitiendo una trazabilidad perfecta de la computación.
\end{fisicaconexion}

\subsection{Transparencia Referencial e Inmutabilidad}
\label{subsec:4_1_2_transparencia}

La transparencia referencial es la propiedad matemática que garantiza que el valor de una expresión depende exclusivamente de sus subexpresiones y no de un historial de estados ocultos o del momento del tiempo en que se evalúa \cite{SICP, PFPL}. En términos prácticos, una expresión es referencialmente transparente si puede ser sustituida por su valor resultante sin alterar el comportamiento del programa \cite{PFPL}.

\begin{definicion}{Igualdad por Sustitución}
Sea una función $\eqfn{f}$ tal que para un argumento $\equ{a}$ produce un valor $\eqk{v}$. En un lenguaje funcional, la evaluación es consistente si para cualquier contexto computacional $\equ{C}$, la siguiente equivalencia se mantiene \cite{PFPL}:
\begin{equation}
    \equ{C}\eqop{[}\eqfn{f}(\equ{a})\eqop{]} \eqop{\equiv} \equ{C}\eqop{[}\eqk{v}\eqop{]}
\end{equation}
donde $\eqop{\equiv}$ denota la identidad de comportamiento. Esto permite realizar optimizaciones como la memorización (\textit{memoization}) o la evaluación perezosa (\textit{lazy evaluation}) con total seguridad \cite{SICP, PFPL}.
\end{definicion}

Esta propiedad está intrínsecamente ligada a la **Inmutabilidad**. En el modelo funcional, no existe el concepto de «celda de memoria» que cambia de valor; en su lugar, se utilizan enlaces (\textit{bindings}) permanentes entre nombres y valores \cite{SICP}. Una vez que se define $\equ{x} \eqop{=} \eqk{13}$, cualquier intento de «cambiar» $\equ{x}$ es, en realidad, la creación de un nuevo ámbito con un nuevo enlace, dejando el valor original intacto \cite{SICP, PFPL}.

\begin{ejemplo}{Transparencia vs. Estado Mutable}
Consideremos la expresión $\equ{R} \eqop{=} \frac{\eqfn{f}(\equ{a}) \eqop{+} \equ{b}}{\eqfn{f}(\equ{a}) \eqop{-} \equ{c}}$. 
\begin{itemize}
    \item \textbf{En FP:} Podemos evaluar $\eqfn{f}(\equ{a})$ una sola vez, asignarlo a un temporal $\equ{t}$, y calcular $\frac{\equ{t} \eqop{+} \equ{b}}{\equ{t} \eqop{-} \equ{c}}$ obteniendo el mismo resultado \cite{SICP}.
    \item \textbf{En Imperativo:} Si $\eqfn{f}$ modifica una variable global o realiza I/O (efecto secundario), la segunda llamada a $\eqfn{f}(\equ{a})$ podría devolver un valor distinto, haciendo que la optimización por sustitución sea incorrecta \cite{PFPL}.
\end{itemize}
\end{ejemplo}

\subsection*{El Modelo de Sustitución ($\beta$-reducción)}

El motor de ejecución de la programación funcional es el **Modelo de Sustitución**, formalizado en el cálculo lambda como la regla de \textbf{$\beta$-reducción} \cite{SICP, PFPL}. Bajo este modelo, aplicar una función a sus argumentos consiste en sustituir cada ocurrencia de los parámetros formales en el cuerpo de la función por los argumentos reales \cite{SICP, TAPL}.

\begin{equation}
    \eqop{(}\eqfn{\lambda} \equ{x} \eqop{.} \equ{M}\eqop{)} \equ{N} \eqop{\twoheadrightarrow} \equ{M}\eqop{[}\equ{N}\eqop{/}\equ{x}\eqop{]}
\end{equation}

Donde $\equ{M}\eqop{[}\equ{N}\eqop{/}\equ{x}\eqop{]}$ representa el término $\equ{M}$ con todas las instancias libres de $\equ{x}$ reemplazadas por $\equ{N}$ \cite{PFPL, TAPL}. Este proceso de «sustituir iguales por iguales» es la base de la razonabilidad del código funcional \cite{HtDP}.

\begin{fisicaconexion}{Invarianza de identidad y conservación}
La transparencia referencial es el equivalente algorítmico a la invarianza de las leyes físicas bajo traslaciones temporales. En un experimento ideal (función pura), si las condiciones iniciales (argumentos) son idénticas, el resultado debe ser idéntico independientemente de «cuándo» se realice el experimento. La inmutabilidad garantiza que la identidad de un objeto de datos se conserve a lo largo de la evolución del sistema, evitando la entropía informativa que surge cuando los estados de los componentes cambian de forma impredecible y no local.
\end{fisicaconexion}

\subsection{El Cálculo Lambda ($\lambda$-calculus) como Sustrato}
\label{subsec:4_1_3_calculo_lambda}

El fundamento matemático de la programación funcional es el \textbf{Cálculo Lambda ($\lambda$-calculus)}, un sistema formal propuesto por Alonzo Church en la década de 1930 para investigar la naturaleza de la computación mediante la aplicación de funciones \cite{PFPL, TAPL}. A pesar de su simplicidad sintáctica, este sistema es un modelo de computación universal equivalente a la Máquina de Turing, lo que implica que cualquier programa ejecutable en hardware convencional puede expresarse íntegramente como un término lambda \cite{PFPL}.

\begin{definicion}{Gramática de Términos Lambda}
El conjunto de todos los términos lambda posibles, denotado como $\eqdom{\Lambda}$, se construye inductivamente a partir de un conjunto infinito de variables mediante tres reglas fundamentales \cite{PFPL, TAPL}:
\begin{enumerate}
    \item \textbf{Variables:} Un identificador $\equ{x}$ es un término. Representa un parámetro o un valor pendiente de resolución \cite{PFPL}.
    \item \textbf{Abstracción ($\eqop{\lambda}$):} Si $\equ{x}$ es una variable y $\equ{M}$ es un término, entonces $\eqop{(}\eqop{\lambda} \equ{x} \eqop{.} \equ{M}\eqop{)}$ es un término que representa una función anónima donde $\equ{x}$ está ligado en el cuerpo $\equ{M}$ \cite{PFPL, TAPL}.
    \item \textbf{Aplicación:} Si $\equ{M}$ y $\equ{N}$ son términos, entonces $\eqop{(}\equ{M} \equ{N}\eqop{)}$ representa la aplicación de la función $\equ{M}$ sobre el argumento $\equ{N}$ \cite{PFPL, TAPL}.
\end{enumerate}
\end{definicion}

La dinámica del cálculo lambda se rige por la \textbf{$\beta$-reducción}, que formaliza el modelo de sustitución introducido previamente \cite{SICP, PFPL}. Para asegurar la consistencia del sistema, se utiliza la \textbf{$\alpha$-conversión}, que permite renombrar variables ligadas para evitar colisiones de nombres durante la evaluación, garantizando que el significado del programa dependa de la estructura y no de las etiquetas arbitrarias de los parámetros \cite{PFPL}.

\begin{ejemplo}{Codificación de Church (Church Booleans)}
Dado que el cálculo lambda puro carece de tipos de datos primitivos, valores como el booleano «verdadero» se codifican como funciones de selección de orden superior \cite{TAPL, PFPL}:
\begin{equation}
    \eqk{true} \eqop{\equiv} \eqop{\lambda} \equ{x} \eqop{.} \eqop{\lambda} \equ{y} \eqop{.} \equ{x}
\end{equation}
Esta función toma dos argumentos y, por definición, siempre selecciona el primero ($\equ{x}$), actuando como la rama «entonces» de un condicional implícito \cite{TAPL}.
\end{ejemplo}

\begin{observacion}{La Tesis de Church y la Definibilidad}
La potencia del cálculo lambda reside en su capacidad de sintetizar estructuras complejas, como números (\textbf{Numerales de Church}) y recursividad (\textbf{Combinador de Punto Fijo} \eqfn{Y}), utilizando únicamente el mecanismo de aplicación funcional \cite{PFPL, TAPL}. Esto convierte al cálculo lambda en el lenguaje «ensamblador» de la programación funcional, sobre el cual se construyen abstracciones de alto nivel como Scheme o Haskell \cite{HtDP, SICP}.
\end{observacion}

\begin{fisicaconexion}{Invarianza de norma y topología computacional}
La $\alpha$-conversión en el cálculo lambda es análoga a la invarianza de norma (\textit{gauge invariance}) en física de campos. Así como el potencial eléctrico absoluto no tiene significado físico —solo las diferencias de potencial importan—, los nombres específicos de las variables ligadas en un término lambda no alteran su capacidad de cómputo. Lo que persiste es la «conectividad» o topología del término: quién provee la información (abstracción) y quién la consume (variable ligada). En este sentido, la computación funcional es una propiedad geométrica y estructural del flujo de datos.
\end{fisicaconexion}

\subsection{Funciones como Ciudadanos de Primera Clase}
\label{subsec:4_1_4_ciudadanos_primera_clase}

La potencia expresiva de \eqk{Scheme} y del paradigma funcional en general emana del estatus otorgado a los procedimientos. Se dice que las funciones son \textbf{ciudadanos de primera clase} (\textit{first-class citizens}) cuando el lenguaje las trata con la misma flexibilidad que a los números o las cadenas de caracteres \cite{SICP, PFPL}.

\begin{definicion}{Propiedades de Primera Clase}
Un elemento de programación posee estatus de primera clase si cumple con los siguientes requisitos funcionales \cite{SICP}:
\begin{enumerate}
    \item \textbf{Asignación:} Puede ser vinculado a un nombre o variable: $\equ{f} \eqop{:=} \eqfn{lambda}(\equ{x}) \dots$
    \item \textbf{Pasaje como Argumento:} Puede ser enviado como parámetro a otro procedimiento (Funciones de Orden Superior) \cite{SICP}.
    \item \textbf{Retorno dinámico:} Puede ser el resultado de la ejecución de una función, permitiendo la creación de clausuras (\textit{closures}) \cite{TAPL, PFPL}.
    \item \textbf{Inclusión en Estructuras:} Puede ser almacenado en listas, árboles o tablas de símbolos \cite{SICP}.
\end{enumerate}
\end{definicion}

Esta capacidad permite lo que se denomina \textbf{abstracción procedimental}. En lugar de escribir múltiples funciones que comparten un patrón similar (como sumar una serie, calcular el producto de un rango o integrar una función), el programador puede diseñar un «meta-procedimiento» que acepte la lógica específica como un argumento funcional \cite{SICP}.

\begin{ejemplo}{Abstracción de Sumatorias (Sigma Notation)}
Siguiendo la lógica de SICP, podemos definir un procedimiento general $\eqfn{sum}$ que encapsula el concepto matemático de sumatoria $\sum_{n=a}^b f(n)$, abstrayendo el término a sumar y el paso de incremento:
\begin{lstlisting}[language=Lisp, caption={Sumatoria como función de orden superior}]
(define (sum term a next b)
  (if (> a b)
      0
      (+ (term a)
         (sum term (next a) next b))))
\end{lstlisting}
Aquí, $\equ{term}$ y $\equ{next}$ no son valores estáticos, sino funciones que el $\eqfn{sum}$ aplicará dinámicamente en cada paso de la recursión \cite{SICP}.
\end{ejemplo}

\begin{observacion}{Hacia los Lenguajes de Dominio Específico (DSL)}
El hecho de que las funciones sean datos permite que el programador extienda el lenguaje base. En la práctica doctoral, esto se traduce en la capacidad de crear lenguajes de dominio específico (\eqk{DSLs}) donde las operaciones complejas se componen de manera modular, facilitando el modelado de sistemas que serían prohibitivamente difíciles de expresar en lenguajes imperativos rígidos \cite{SICP, HtDP}.
\end{observacion}

\begin{fisicaconexion}{Operadores funcionales y mecánica cuántica}
En mecánica cuántica, los observables físicos (como el momento o la energía) no se representan como números, sino como operadores lineales que actúan sobre funciones de onda. Al igual que en la programación funcional, estos operadores son «ciudadanos de primera clase» de la teoría: pueden sumarse, multiplicarse (componerse) y sus propiedades se derivan de su estructura algebraica. El \textit{Hamiltoniano} ($\eqfn{H}$), por ejemplo, es una función de orden superior que determina la evolución temporal del sistema; su aplicación a un estado produce la derivada temporal de dicho estado, ilustrando una simetría profunda entre el procesamiento de información y las leyes de la naturaleza.
\end{fisicaconexion}

\section{Introducción a Scheme y DrRacket}
\label{sec:4_2_intro_scheme}

Scheme es un dialecto de \eqk{LISP} surgido a mediados de la década de 1970, diseñado para ser un lenguaje funcional minimalista con una semántica extraordinariamente limpia \cite{SICP}. A diferencia de otros lenguajes, Scheme reduce la computación a un conjunto pequeño de reglas de combinación, permitiendo que estructuras complejas emerjan de la composición de elementos simples \cite{SICP, HtDP}.

\subsection{El Entorno de Programación DrRacket}
\label{subsec:4_2_1_drracket}

El vehículo moderno para el estudio de Scheme es \eqk{DrRacket}, un entorno de desarrollo integrado que no solo proporciona un editor, sino un ecosistema de lenguajes graduados diseñados para el aprendizaje sistemático del diseño de programas \cite{HtDP}.

Una de las decisiones de diseño fundamentales en Scheme es el compromiso con el \textbf{tipado dinámico} \cite{SICP}:

\begin{definicion}{Tipado Dinámico}
En Scheme, los tipos están asociados a los \textbf{valores} y no a las variables o nombres \cite{SICP}. Esto implica que el sistema de tipos no realiza una verificación estática en tiempo de compilación; en su lugar, la seguridad de tipos se garantiza mediante comprobaciones durante la ejecución (\textit{runtime checks}), lo que otorga una flexibilidad expresiva superior a costa de postergar la detección de errores de tipo hasta el momento en que se procesan los datos \cite{SICP, HtDP}.
\end{definicion}

El flujo de trabajo en DrRacket se basa en el ciclo \eqk{REPL} (\textit{Read-Eval-Print Loop}), el cual actúa como un motor de transformación de estado informativo (ver Figura \ref{fig:repl_flow}).

\begin{figure}[htbp]
    \centering
    \begin{tikzpicture}[node distance=2.5cm, >=latex, font=\small]
        % Nodos del flujo
        \node[draw, circle, fill=CtpBase] (read) {Read};
        \node[draw, circle, fill=CtpMantle, right of=read] (eval) {Eval};
        \node[draw, circle, fill=CtpBase, right of=eval] (print) {Print};
        
        % Conexiones
        \draw[->] (read) -- node[above] {S-expr} (eval);
        \draw[->] (eval) -- node[above] {Valor} (print);
        \draw[->] (print) .. controls +(0,-1.5) and +(0,-1.5) .. node[below] {Esperar entrada} (read);
        
        % Usuario
        \node[left of=read, xshift=-1cm] (user) {Usuario};
        \draw[->, dashed] (user) -- (read);
    \end{tikzpicture}
    \caption{Ciclo REPL: El motor de interacción fundamental de Scheme \cite{SICP, HtDP}.}
    \label{fig:repl_flow}
\end{figure}

\begin{observacion}{Eficiencia y Gestión de Memoria}
A pesar de su naturaleza dinámica, las implementaciones modernas de Scheme son altamente eficientes gracias a la compilación a código nativo y a la gestión automática de memoria mediante \textbf{recolección de basura} (\textit{garbage collection}) \cite{SICP, HtDP}. El sistema detecta automáticamente los objetos que ya no son accesibles desde el entorno global o la pila de llamadas, recuperando el espacio de memoria sin intervención explícita del programador \cite{SICP}.
\end{observacion}

\subsection{La Anatomía de una S-expresión}
\label{subsec:4_2_2_sexpresion}

La sintaxis de \eqk{Scheme} es notablemente minimalista y uniforme, basándose casi exclusivamente en el concepto de \textbf{S-expresión} \cite{SICP}. Este diseño permite que el lenguaje sea homoicónico, es decir, que la representación del código sea idéntica a la estructura de datos que el propio lenguaje manipula \cite{SICP, HtDP}.

\begin{definicion}{S-expresión (Symbolic Expression)}
Una S-expresión es una estructura definida de forma inductiva como \cite{SICP}:
\begin{enumerate}
    \item \textbf{Átomo:} Un elemento indivisible, como un número ($\eqk{42}$), un booleano ($\eqk{\#t}$), una cadena de texto o un símbolo.
    \item \textbf{Lista:} Una secuencia de S-expresiones encerrada entre paréntesis $\eqop{(} \dots \eqop{)}$, donde los elementos están separados por espacios en blanco \cite{SICP}.
\end{enumerate}
\end{definicion}

A diferencia de la notación infija convencional utilizada en matemáticas y lenguajes imperativos (ej. $3 + 4$), Scheme utiliza la \textbf{notación prefija} o de Cambridge \cite{SICP, HtDP}. En esta convención, el primer elemento de una lista se interpreta como el operador o función a aplicar, y los elementos restantes como sus argumentos.

\begin{ejemplo}{Composición y Notación Prefija}
Consideremos la expresión aritmética compleja $3 + 4 \times (5 - 2)$. En Scheme, esta jerarquía se expresa mediante el anidamiento de S-expresiones \cite{SICP, HtDP}:
\begin{equation}
    \eqop{(}\eqfn{+} \eqk{3} \eqop{(}\eqfn{*} \eqk{4} \eqop{(}\eqfn{-} \eqk{5} \eqk{2}\eqop{)}\eqop{)}\eqop{)}
\end{equation}
El intérprete evalúa esta estructura de forma recursiva: primero resuelve las subexpresiones más internas y luego aplica los operadores externos a los resultados obtenidos \cite{SICP}.
\end{ejemplo}

\begin{observacion}{Regla de Evaluación de Combinaciones}
Para evaluar una combinación (una lista no vacía), el intérprete sigue un procedimiento recursivo \cite{SICP}:
\begin{enumerate}
    \item Evalúa cada una de las subexpresiones de la combinación.
    \item Aplica el procedimiento (valor de la subexpresión más a la izquierda) a los argumentos (valores de las demás subexpresiones).
\end{enumerate}
Esta naturaleza recursiva permite que Scheme maneje niveles arbitrarios de anidamiento sin necesidad de reglas de precedencia de operadores complejas, ya que los paréntesis dictan explícitamente el orden de evaluación \cite{SICP, HtDP}.
\end{observacion}

\begin{fisicaconexion}{Jerarquías de escala y fractales sintácticos}
La estructura de las S-expresiones posee una propiedad de auto-similitud análoga a la de los fractales en física. Un fractal es una figura geométrica donde cada parte, por pequeña que sea, conserva la misma estructura que el todo. De forma similar, en una S-expresión, cada sub-lista sigue exactamente las mismas reglas gramaticales y de evaluación que la expresión raíz. Esta invarianza de escala sintáctica simplifica enormemente el motor de inferencia (el evaluador) y permite construir abstracciones poderosas mediante la composición recursiva de bloques elementales.
\end{fisicaconexion}

\subsection{Enlaces y Definiciones Globales}
\label{subsec:4_2_3_enlaces}

En el paradigma funcional, la abstracción se logra mediante la asociación de nombres descriptivos con los resultados de expresiones complejas, permitiendo al programador referirse a un objeto computacional sin necesidad de recalcularlo o conocer su estructura interna \cite{SICP}. En \eqk{Scheme}, este proceso se realiza a través de la forma especial \eqop{define}, la cual crea un enlace (\textit{binding}) en el entorno global \cite{SICP, HtDP}.

\begin{definicion}{Definición de Constantes}
La sintaxis general para vincular un nombre a un valor es:
\begin{equation}
    \eqop{(}\eqop{define} \equ{nombre} \equ{expresion}\eqop{)}
\end{equation}
A diferencia de la asignación imperativa, esta operación no representa una modificación de una celda de memoria existente, sino la creación de una asociación permanente e inmutable entre el identificador $\equ{nombre}$ y el valor resultante de evaluar $\equ{expresion}$ \cite{SICP, HtDP}.
\end{definicion}

El uso de constantes es una práctica recomendada en el diseño sistemático de programas para evitar la duplicación de «números mágicos» y facilitar el mantenimiento del código \cite{HtDP}. 

\begin{ejemplo}{Enlaces globales y legibilidad}
Consideremos la definición de parámetros físicos para una simulación \cite{HtDP}:
\begin{lstlisting}[language=Lisp, caption={Definiciones de constantes en Scheme}]
(define pi 3.14159)
(define radio 10)
(define area (* pi (* radio radio))) ; El valor se calcula una sola vez
\end{lstlisting}
En este fragmento, el intérprete liga \equ{pi} y \equ{radio} a sus respectivos literales, y luego evalúa la expresión del área utilizando los enlaces previos \cite{SICP}.
\end{ejemplo}

\subsection*{Abstracción Procedural (Definición de Funciones)}

Aunque el cálculo lambda puro define funciones exclusivamente mediante la abstracción \eqop{lambda}, \eqk{Scheme} proporciona una notación abreviada —comúnmente denominada azúcar sintáctico— que permite definir procedimientos de una manera más cercana a la notación matemática tradicional \cite{SICP}.

\begin{definicion}{Sintaxis de Definición de Funciones}
La forma abreviada $\eqop{(}\eqop{define} \eqop{(}\equ{f} \equ{x_1} \dots \equ{x_n}\eqop{)} \equ{cuerpo}\eqop{)}$ es semánticamente equivalente a realizar un enlace de una función anónima a un nombre \cite{SICP}:
\begin{equation}
\eqop{define} \equ{f} \eqop{(}\eqfn{lambda} \eqop{(}\equ{x_1} \dots \equ{x_n}\eqop{)} \equ{cuerpo}\eqop{)}
\end{equation}
Esta uniformidad garantiza que las funciones sean tratadas como datos y que el modelo de evaluación por sustitución se mantenga consistente \cite{SICP}.
\end{definicion}

\begin{observacion}{Predicados de Tipo}
Como parte de su biblioteca estándar, Scheme incluye funciones especializadas llamadas \textbf{predicados}, que devuelven un valor booleano ($\eqk{\#t}$ o $\eqk{\#f}$) al evaluar una propiedad de un dato \cite{HtDP, SICP}. Ejemplos comunes incluyen \eqfn{number?}, \eqfn{zero?}, \eqfn{equal?} y \eqfn{null?}, los cuales son esenciales para implementar la lógica condicional en programas funcionales \cite{SICP}.
\end{observacion}

\begin{fisicaconexion}{Identidad de partículas y etiquetas de campo}
El mecanismo de \eqop{define} en programación funcional guarda una analogía profunda con la asignación de etiquetas a campos escalares en física teórica. Cuando definimos un valor como \equ{pi} o una función como \eqfn{f}, no estamos creando una «caja» que puede cambiar su contenido (como en el modelo imperativo), sino que estamos nombrando una propiedad intrínseca del universo del programa. Al igual que una constante universal (como la constante de Planck $\eqcon{h}$), el nombre es simplemente una referencia a una verdad inmutable. En este sentido, un programa funcional no es una historia de cambios, sino un conjunto de identidades estáticas que interactúan, similar a un sistema de ecuaciones de campo donde las variables representan estados fijos en el espacio-tiempo.
\end{fisicaconexion}

\subsection{Control de Flujo como Expresión: \texttt{if} y \texttt{cond}}
\label{subsec:4_2_4_control_flujo}

A diferencia de los lenguajes imperativos, donde el control de flujo se basa en «sentencias» que dirigen el puntero de instrucción sin retornar necesariamente un valor, en \eqk{Scheme} todas las estructuras condicionales son **expresiones** \cite{HtDP, SICP}. Esto implica que un condicional puede ser utilizado en cualquier lugar donde se espere un dato, integrándose naturalmente en la composición de funciones \cite{HtDP}.

\begin{definicion}{La Expresión \texttt{if}}
La forma \eqop{(}\eqop{if} \equ{predicado} \equ{consecuente} \equ{alternativa}\eqop{)} es la estructura de decisión binaria fundamental \cite{SICP}. Su evaluación ocurre en dos etapas estrictas:
\begin{enumerate}
    \item Se evalúa el \equ{predicado}, el cual debe producir un valor booleano (\eqk{\#t} o \eqk{\#f}).
    \item Si el resultado es \eqk{\#t}, se evalúa únicamente el \equ{consecuente}. Si es \eqk{\#f}, se evalúa la \equ{alternativa}. El valor resultante se convierte en el valor de toda la expresión \eqop{if} \cite{SICP}.
\end{enumerate}
\end{definicion}

Para escenarios que requieren múltiples ramificaciones basadas en distintas clases de datos, Scheme proporciona la forma especial \eqop{cond} \cite{HtDP}.

\begin{definicion}{La Expresión \texttt{cond}}
La forma \eqop{cond} permite una selección múltiple mediante una secuencia de cláusulas \cite{SICP, HtDP}:
\[ \eqop{(}\eqop{cond} \eqop{[}\equ{P_1} \equ{E_1}\eqop{]} \dots \eqop{[}\eqop{else} \equ{E_n}\eqop{]}\eqop{)} \]
El intérprete evalúa cada predicado $\equ{P_i}$ en orden. Al encontrar el primero que sea \eqk{\#t}, ejecuta su correspondiente expresión $\equ{E_i}$ y termina \cite{SICP}. La cláusula \eqop{else} es una salvaguarda que garantiza que la expresión siempre retorne un valor si ninguna condición previa se cumple \cite{HtDP, SICP}.
\end{definicion}

\begin{ejemplo}{Lógica de semáforos mediante \texttt{cond}}
Siguiendo el diseño sistemático de \cite{HtDP}, podemos definir la transición de estados de un sistema mediante \eqop{cond}, donde cada rama corresponde a un caso del dominio de datos \cite{HtDP}:
\begin{lstlisting}[language=Lisp, caption={Transición de estados en Scheme}]
(define (next-state color)
  (cond
    [(string=? color "red") "green"]
    [(string=? color "green") "yellow"]
    [(string=? color "yellow") "red"]
    [else (error "Estado invalido")]))
\end{lstlisting}
\end{ejemplo}

\begin{observacion}{Pragmatismo: \texttt{if} vs. \texttt{cond}}
Aunque \eqop{if} puede anidarse para simular múltiples ramas, se considera una mejor práctica de ingeniería utilizar \eqop{if} solo para decisiones de tipo «uno u otro» y reservar \eqop{cond} para casos donde la lógica emana de definiciones de datos complejas o múltiples intervalos \cite{HtDP, SICP}.
\end{observacion}

\begin{fisicaconexion}{Sistemas de conmutación y potenciales de paso}
El control de flujo en programación funcional guarda una analogía directa con los potenciales de paso en mecánica cuántica o los sistemas de conmutación electrónicos. En un pozo de potencial, el «camino» o estado final de una partícula depende estrictamente de si su energía supera un umbral crítico (el predicado). Al igual que en Scheme, donde el resultado es una función directa de la entrada sin estados ocultos, en estos sistemas físicos la trayectoria se bifurca determinísticamente basándose en las condiciones de contorno iniciales, sin «memoria» del proceso de decisión previo.
\end{fisicaconexion}


\section{Procesamiento de Listas y Datos Estructurados}
\label{sec:4_3_listas}

En el paradigma funcional, la lista es la estructura de datos primaria, valorada por su flexibilidad y su naturaleza inherentemente recursiva \cite{SICP, HtDP}. A diferencia de los arreglos en lenguajes imperativos, las listas en \eqk{Scheme} no ocupan bloques de memoria contiguos, sino que se construyen mediante el encadenamiento de \textbf{celdas de par} (\textit{cons cells}), permitiendo una manipulación dinámica eficiente sin necesidad de redimensionamiento explícito \cite{SICP}.

\subsection{Primitivas de Construcción y Deconstrucción: \texttt{cons}, \texttt{car} y \texttt{cdr}}
\label{subsec:4_3_1_primitivas}

Toda la manipulación de datos estructurados en Scheme se reduce a tres operaciones fundamentales que actúan sobre pares \cite{SICP}:

\begin{itemize}
    \item \textbf{\eqfn{cons} (\textit{Construct}):} Toma dos argumentos y devuelve un nuevo par. Es la unidad básica de «pegamento» de datos \cite{SICP}.
    \item \textbf{\eqfn{car} (\textit{Contents of Address Register}):} Extrae el primer elemento de un par \cite{SICP}.
    \item \textbf{\eqfn{cdr} (\textit{Contents of Data Register}):} Extrae el segundo elemento del par (típicamente el resto de la lista) \cite{SICP}.
\end{itemize}

\begin{ejemplo}{Construcción manual de listas}
Para representar la secuencia $\eqk{1, 2, 3}$, debemos anidar las operaciones de construcción hasta alcanzar el delimitador de lista vacía \eqk{'()} \cite{SICP}:
\begin{equation}
    \equ{L} \eqop{:=} \eqop{(}\eqfn{cons} \eqk{1} \eqop{(}\eqfn{cons} \eqk{2} \eqop{(}\eqfn{cons} \eqk{3} \eqk{'()}\eqop{)}\eqop{)}\eqop{)}
\end{equation}
Al aplicar los selectores, obtenemos:
\begin{itemize}
    \item $\eqop{(}\eqfn{car} \equ{L}\eqop{)} \eqop{\rightarrow} \eqk{1}$
    \item $\eqop{(}\eqfn{cdr} \equ{L}\eqop{)} \eqop{\rightarrow} \eqop{(}\eqk{2 3}\eqop{)}$
    \item $\eqop{(}\eqfn{car} \eqop{(}\eqfn{cdr} \equ{L}\eqop{)}\eqop{)} \eqop{\rightarrow} \eqk{2}$
\end{itemize}
\end{ejemplo}

\begin{figure}[htbp]
    \centering
    \begin{tikzpicture}[list/.style={rectangle split, rectangle split parts=2, draw, rectangle split horizontal, fill=CtpBase}, >=latex]
        % Celda 1
        \node[list] (cell1) {$\eqk{1}$ \nodepart{two}};
        % Celda 2
        \node[list, right of=cell1, xshift=1.5cm] (cell2) {$\eqk{2}$ \nodepart{two}};
        % Celda 3
        \node[list, right of=cell2, xshift=1.5cm] (cell3) {$\eqk{3}$ \nodepart{two}};
        % Null
        \node[draw, circle, right of=cell3, xshift=1cm] (nil) {$\eqk{\emptyset}$};

        % Conexiones
        \draw[->] (cell1.second east) -- (cell2.west);
        \draw[->] (cell2.second east) -- (cell3.west);
        \draw[->] (cell3.second east) -- (nil);
        
        \node[above of=cell1, yshift=-0.5cm] {car};
        \node[right of=cell1, xshift=0.5cm, yshift=0.3cm] {cdr};
    \end{tikzpicture}
    \caption{Representación de «Caja y Puntero» de una lista en Scheme \cite{SICP}.}
    \label{fig:box_and_pointer}
\end{figure}

\begin{observacion}{Homoiconicidad y el operador \eqop{quote}}
Scheme permite tratar el código como datos mediante el operador de cita $\eqop{'}$ \cite{SICP}. Mientras que $\eqop{(}\eqfn{+} \eqk{1 2}\eqop{)}$ evalúa a $\eqk{3}$, la expresión $\eqop{'}\eqop{(}\eqfn{+} \eqk{1 2}\eqop{)}$ devuelve una lista literal que contiene el símbolo \eqfn{+} y los números \eqk{1} y \eqk{2} \cite{SICP}. Esta propiedad es fundamental para la meta-programación y la creación de macros.
\end{observacion}


\subsection{Estructura Inductiva de las Listas}
\label{subsec:4_3_2_estructura_inductiva}

La potencia de las listas en la programación funcional no reside únicamente en su implementación física, sino en su \textbf{definición inductiva}. Bajo este prisma, una lista no es un contenedor estático, sino un objeto matemático construido mediante reglas de inferencia claras \cite{PFPL, HtDP}.

\begin{definicion}{Definición Inductiva de Lista}
El conjunto de todas las listas, denotado como $\eqdom{List}$, es el juicio más pequeño cerrado bajo las siguientes reglas de formación \cite{PFPL, HtDP}:
\begin{enumerate}
    \item \textbf{Base:} La lista vacía $\eqk{'()}$ (o \textit{nil}) pertenece a $\eqdom{List}$.
    \item \textbf{Inducción:} Si $\equ{e}$ es un átomo y $\equ{L}$ es una lista ($\equ{L} \eqop{\in} \eqdom{List}$), entonces $\eqop{(}\eqfn{cons} \equ{e} \equ{L}\eqop{)}$ también pertenece a $\eqdom{List}$.
\end{enumerate}
\end{definicion}

Esta naturaleza inductiva otorga a las listas la propiedad de \textbf{clausura}: los resultados de combinar objetos de datos pueden ser, a su vez, combinados utilizando el mismo mecanismo de construcción \cite{SICP}. Como resultado, Scheme permite la creación de estructuras jerárquicas complejas, como listas de listas o árboles, manteniendo una uniformidad semántica total \cite{SICP, HtDP}.

\begin{observacion}{Eficiencia de Acceso}
Debido a esta estructura de celdas ligadas, el acceso al primer elemento ($\eqfn{car}$) es una operación de tiempo constante $\eqop{O(1)}$. Sin embargo, acceder al $\equ{n}$-ésimo elemento requiere un recorrido lineal $\eqop{O(n)}$, lo que diferencia fundamentalmente a las listas de los arreglos imperativos de acceso aleatorio \cite{SICP}.
\end{observacion}

\subsection{El Patrón de Recursión Estructural}
\label{subsec:4_3_3_recursion_estructural}

La estructura de un algoritmo funcional suele ser una imagen especular de la estructura de los datos que procesa \cite{HtDP}. Dado que las listas se definen inductivamente, la técnica primordial para manipularlas es la \textbf{recursión estructural}, la cual garantiza que el programa cubra exhaustivamente todos los casos posibles de la definición del dato \cite{HtDP}.

\begin{definicion}{Template de Procesamiento de Listas}
Siguiendo la metodología de diseño sistemático, una función que procesa una lista $\equ{L}$ debe seguir un esquema de análisis de casos basado en predicados estructurales \cite{HtDP}:
\begin{lstlisting}[language=Lisp, caption={Esquema general de recursión estructural}]
(define (process-list L)
  (cond
    [(null? L) ...]                ; Caso Base: lista vacia
    [else                          ; Caso Inductivo: lista no vacia
     (... (car L)                  ; 1. Procesar la cabeza
          (process-list (cdr L)))] ; 2. Recurrir sobre el resto
  )
)
\end{lstlisting}
\end{definicion}

Este patrón asegura que el problema complejo se reduzca en cada paso, acercándose sistemáticamente al caso base $\eqfn{null?}$, lo que en la teoría de computación es fundamental para demostrar la \textbf{terminación} del algoritmo \cite{PFPL}.

\begin{ejemplo}{Cálculo de la longitud de una lista}
Para implementar \eqfn{length}, aplicamos el patrón estructural: si la lista es vacía, su tamaño es $\eqk{0}$; de lo contrario, el tamaño es $\eqk{1}$ más el tamaño del resto de la lista \cite{SICP}:
\begin{lstlisting}[language=Lisp]
(define (length L)
  (if (null? L)
      0
      (+ 1 (length (cdr L)))))
\end{lstlisting}
\end{ejemplo}

\begin{fisicaconexion}{Decantación y filtrado por etapas}
El patrón de recursión estructural es análogo a un proceso de filtrado por etapas en ingeniería química o mecánica de fluidos. La entrada (la lista $\equ{L}$) se somete a una membrana (el análisis de casos); el elemento retenido ($\eqfn{car}$) se procesa inmediatamente, mientras que el resto del fluido ($\eqfn{cdr}$) pasa a la siguiente etapa idéntica de filtración. El proceso concluye cuando no queda más material que filtrar (el caso base). Así como la eficiencia de un filtro depende de la uniformidad de sus etapas, la elegancia de un programa funcional emana de la aplicación repetida de una única ley de transformación sobre una estructura auto-similar.
\end{fisicaconexion}

\subsection{Operaciones Fundamentales y Composición}
\label{subsec:4_3_4_operaciones}

Una vez establecido el patrón de recursión estructural, podemos definir las operaciones que constituyen el núcleo del procesamiento de listas en \eqk{Scheme}. Estas funciones no solo son herramientas de utilidad, sino que ejemplifican cómo la computación funcional transforma estructuras de datos mediante la aplicación sucesiva de reglas de combinación puras \cite{SICP, HtDP}.

\begin{definicion}{Concatenación (\texttt{append})}
La función \eqfn{append} toma dos listas, $\equ{L_1}$ y $\equ{L_2}$, y devuelve una nueva lista que contiene los elementos de $\equ{L_1}$ seguidos de los de $\equ{L_2}$ \cite{SICP}. Su lógica consiste en recorrer $\equ{L_1}$ hasta el final y sustituir su delimitador nulo por $\equ{L_2}$ \cite{SICP}:
\begin{lstlisting}[language=Lisp]
(define (append L1 L2)
  (if (null? L1)
      L2
      (cons (car L1) (append (cdr L1) L2))))
\end{lstlisting}
\end{definicion}

\begin{definicion}{Pertenencia (\texttt{member})}
La función \eqfn{member} verifica si un átomo dado $\equ{atm}$ se encuentra dentro de una lista $\equ{lis}$ \cite{SICP}. Retorna \eqk{\#t} si existe una coincidencia según el predicado \eqfn{equal?} y \eqk{\#f} en caso contrario:
\begin{lstlisting}[language=Lisp]
(define (member atm lis)
  (cond
    [(null? lis) #f]
    [(equal? atm (car lis)) #t]
    [else (member atm (cdr lis))]))
\end{lstlisting}
\end{definicion}

\subsection*{Composición: Lógica de Conjuntos}

La verdadera potencia del paradigma funcional reside en la capacidad de componer funciones simples para construir abstracciones de mayor nivel \cite{SICP, HtDP}. Un ejemplo clásico es la implementación de la \textbf{Diferencia de Conjuntos} (\eqfn{setdiff}), la cual utiliza \eqfn{member} como sub-procedimiento para filtrar elementos \cite{HtDP}.

\begin{ejemplo}{Diferencia de listas mediante filtrado recursivo}
Dadas dos listas, \eqfn{setdiff} construye una nueva lista que contiene solo los elementos de la primera que no están presentes en la segunda \cite{HtDP}:
\begin{lstlisting}[language=Lisp, caption={Implementación de setdiff en Scheme}]
(define (setdiff lis1 lis2)
  (cond
    [(null? lis1) '()]
    [(null? lis2) lis1]
    [(member (car lis1) lis2) (setdiff (cdr lis1) lis2)]
    [else (cons (car lis1) (setdiff (cdr lis1) lis2))]))
\end{lstlisting}
Este procedimiento demuestra cómo la lógica de decisión (\eqop{cond}) dirige el flujo de datos hacia la construcción (\eqfn{cons}) o la omisión de elementos basándose en predicados de orden superior \cite{SICP}.
\end{ejemplo}

\begin{ejemplo}{Traza de evaluación por sustitución}
Para la expresión \texttt{(member 'b '(a b))}, la traza de evaluación ilustra el despliegue de la recursión \cite{SICP}:
\begin{itemize}
    \item \texttt{(member 'b '(a b))}
    \item $\eqop{\rightarrow}$ \texttt{(cond [\#f ...] [(equal? 'b 'a) ...] [else (member 'b '(b))])}
    \item $\eqop{\rightarrow}$ \texttt{(member 'b '(b))}
    \item $\eqop{\rightarrow}$ \texttt{(cond [\#f ...] [(equal? 'b 'b) \#t] ...)}
    \item $\eqop{\rightarrow}$ \eqk{\#t}
\end{itemize}
\end{ejemplo}

\begin{fisicaconexion}{Procesos en serie y decantación diferencial}
La composición de funciones como \eqfn{member} y \eqfn{setdiff} es análoga a una torre de destilación o un sistema de procesos en serie en ingeniería de sistemas. Cada función actúa como un «estadio» de procesamiento con una selectividad específica. En la decantación diferencial, una mezcla pasa por múltiples filtros; cada etapa separa componentes basados en propiedades físicas (masa, tamaño), de forma idéntica a como \eqfn{setdiff} separa átomos basados en el predicado de igualdad. La inmutabilidad de la entrada garantiza que el fluido original no se contamine, permitiendo que el mismo flujo de datos sea procesado por filtros paralelos sin interferencias, una propiedad esencial para el diseño de reactores químicos y sistemas paralelos masivos.
\end{fisicaconexion}



\section{Abstracción de Orden Superior}
\label{sec:4_4_orden_superior}

El §\ref{subsec:4_1_4_ciudadanos_primera_clase} estableció formalmente que las funciones son ciudadanos de primera clase en \eqk{Scheme}: pueden asignarse a nombres, pasarse como argumentos y retornarse como resultados. Esta sección materializa esa promesa operacionalmente. Si las funciones son datos, existe toda una maquinaria de abstracción que opera \textit{sobre} funciones, permitiendo capturar patrones computacionales recurrentes con una generalidad que resultaría inalcanzable en el paradigma imperativo \cite{SICP, HtDP}.

\subsection{Funciones Anónimas: \texttt{lambda} y el Modelo de Clausura}
\label{subsec:4_4_1_lambda}

En §\ref{subsec:4_1_3_calculo_lambda} se introdujo el término de abstracción del $\lambda$-calculus como la construcción $\eqop{(}\eqfn{\lambda}\ \equ{x} \eqop{.} \equ{M}\eqop{)}$, que describe una función anónima que liga $\equ{x}$ en el cuerpo $\equ{M}$. En \eqk{Scheme}, esta construcción no es una metáfora pedagógica: la forma especial \eqop{lambda} \textit{es} la realización directa y ejecutable de ese término \cite{SICP, PFPL}.

\begin{definicion}{Expresión \texttt{lambda} en Scheme}
La forma $\eqop{(}\eqfn{lambda}\ \eqop{(}\equ{x_1} \dots \equ{x_n}\eqop{)}\ \equ{cuerpo}\eqop{)}$ construye un objeto procedural anónimo que, cuando se aplica a $n$ argumentos $\eqk{v_1} \dots \eqk{v_n}$, evalúa $\equ{cuerpo}$ en un nuevo entorno donde cada $\equ{x_i}$ está ligado a $\eqk{v_i}$ \cite{SICP}. A diferencia de \eqop{define}, que crea un enlace en el entorno global, \eqop{lambda} produce un procedimiento como valor de primera clase sin asignarle nombre alguno.
\end{definicion}

La propiedad que eleva a \eqop{lambda} más allá de una mera función anónima es la \textbf{clausura} (\textit{closure}). Cuando el cuerpo de una expresión \eqop{lambda} contiene variables que no son parámetros formales —denominadas \textbf{variables libres}— el intérprete no las resuelve en el momento de la definición, sino que captura el \textit{entorno léxico circundante} como parte del objeto procedural resultante \cite{SICP, TAPL}. El resultado es un par $\langle \text{código},\ \text{entorno} \rangle$ que retiene acceso permanente a los valores que existían en el ámbito de su creación, incluso si ese ámbito ya no es alcanzable desde el programa principal.

\begin{ejemplo}{Generador de funciones mediante clausura}
El procedimiento \eqfn{make-adder} ilustra la captura de entorno: retorna una \eqop{lambda} cuya variable libre \equ{n} queda ligada al argumento de la invocación exterior \cite{SICP}:
\begin{lstlisting}[language=Lisp, caption={Clausura: captura del entorno léxico}]
(define (make-adder n)
  (lambda (x) (+ x n)))

(define add5  (make-adder 5))
(define add10 (make-adder 10))

(add5  3)   ; -> 8   (n = 5 capturado)
(add10 3)   ; -> 13  (n = 10 capturado)
\end{lstlisting}
Cada llamada a \eqfn{make-adder} produce una clausura distinta: \eqfn{add5} y \eqfn{add10} comparten el mismo código pero poseen entornos independientes donde \equ{n} toma valores distintos. La inmutabilidad de \equ{n} garantiza que las clausuras son funciones puras en el sentido del §\ref{subsec:4_1_1_naturaleza} \cite{SICP}.
\end{ejemplo}

La clausura es también la base de la \textbf{aplicación parcial}: dado un procedimiento de $n$ argumentos, es posible fijar los primeros $k < n$ argumentos mediante una clausura que capture esos valores, obteniendo un procedimiento de $n - k$ argumentos. Este patrón —conocido como \textit{currying} en honor a Haskell Curry— es ubicuo en lenguajes funcionales y permite construir familias de funciones especializadas a partir de una única definición general \cite{SICP, PFPL}.

El mecanismo de resolución de variables libres depende críticamente de la regla de ámbito del lenguaje. \eqk{Scheme} utiliza \textbf{ámbito estático} (\textit{lexical scoping}): el entorno relevante es aquel en el que la expresión \eqop{lambda} fue \textit{definida}, no aquel en el que es \textit{invocada} \cite{SICP}. Esta decisión es lo que hace que las clausuras sean semánticamente predecibles — el significado de una función depende exclusivamente de su posición textual en el código fuente, no del contexto de llamada.

\begin{fisicaconexion}{Estados coherentes y amplitudes fijas de captura}
Una clausura es el análogo computacional de un estado coherente en óptica cuántica. Un estado coherente $|\alpha\rangle$ no es un eigenestado de un observable, sino una superposición de eigenstados $|n\rangle$ con amplitudes $c_n$ determinadas en el momento de su preparación, no en el momento de la medición. De manera análoga, una clausura no es un procedimiento descontextualizado: es el par procedimiento-entorno donde las amplitudes de las variables libres quedan fijadas en el instante de la abstracción. Cualquier invocación posterior mide el procedimiento en ese entorno capturado, con resultados determinados por él, independientemente del contexto de la medición.
\end{fisicaconexion}

\subsection{\texttt{let} como Abstracción de Aplicación Inmediata}
\label{subsec:4_4_2_let}

La forma especial \eqop{let} es \textbf{azúcar sintáctico} sobre \eqop{lambda}: introduce enlaces locales sin crear un procedimiento persistente ni modificar el entorno global \cite{SICP, HtDP}. La equivalencia semántica es exacta:

\begin{equation}
    \eqop{(}\eqk{let}\ \eqop{((}\equ{x_1}\ \eqk{v_1}\eqop{)} \cdots \eqop{(}\equ{x_n}\ \eqk{v_n}\eqop{))}\ \equ{M}\eqop{)} \;\eqop{\equiv}\; \eqop{((}\eqfn{\lambda}\ \eqop{(}\equ{x_1} \cdots \equ{x_n}\eqop{)}\ \equ{M}\eqop{)}\ \eqk{v_1} \cdots \eqk{v_n}\eqop{)}
\end{equation}

El intérprete evalúa todas las expresiones de valor $\eqk{v_1} \dots \eqk{v_n}$ en el entorno \textit{exterior} al \eqop{let} —sin acceso mutuo entre ellas— crea un nuevo marco con los enlaces $\equ{x_i} \mapsto \eqk{v_i}$, evalúa el cuerpo $\equ{M}$ en ese marco, y descarta el marco sin dejar rastro en el entorno global \cite{SICP}. No hay asignación; hay sustitución en el sentido del §\ref{subsec:4_1_2_transparencia}.

\begin{observacion}{Variantes de \texttt{let} y el orden de evaluación}
El orden en que los bindings se hacen visibles da lugar a tres variantes con semánticas distintas \cite{SICP, HtDP}:
\begin{itemize}
    \item \eqk{let}: todos los $\eqk{v_i}$ se evalúan en el entorno exterior; los bindings son mutuamente invisibles. Es el caso base descrito arriba.
    \item \eqk{let*}: los bindings se introducen secuencialmente; $\eqk{v_i}$ puede referirse a $\equ{x_j}$ con $j < i$. Equivale a \eqop{let} anidados.
    \item \eqk{letrec}: todos los nombres son visibles en todas las expresiones de valor desde el inicio, permitiendo definiciones locales mutuamente recursivas. Es el mecanismo necesario para funciones auxiliares que se invocan entre sí \cite{SICP}.
\end{itemize}
\end{observacion}

\begin{ejemplo}{Aplicación de orden superior con \texttt{let}}
El siguiente fragmento liga mediante \eqop{let} una \eqop{lambda} de orden superior como valor local, demostrando la composición de ambas formas \cite{SICP}:
\begin{lstlisting}[language=Lisp, caption={let con función de orden superior}]
(let ((double-any (lambda (f x) (f x x))))
  (list (double-any + 13)
        (double-any cons 'a)))
; -> (26 (a . a))
\end{lstlisting}
\eqfn{double-any} recibe una función \equ{f} y un valor \equ{x}, y aplica \equ{f} a \equ{x} consigo mismo. Al recibir \eqop{+} produce $13 + 13 = 26$; al recibir \eqfn{cons} produce el par \eqk{(a . a)}. El mismo procedimiento opera sobre datos de naturaleza completamente distinta gracias al tipado dinámico de \eqk{Scheme} \cite{SICP}.
\end{ejemplo}

\subsection{Patrones de Iteración Funcional: \texttt{map}, \texttt{filter} y \texttt{fold}}
\label{subsec:4_4_3_hof_patrones}

El reconocimiento de que ciertos patrones de recursión estructural se repiten invariablemente en el procesamiento de listas conduce a una de las abstracciones más poderosas del paradigma funcional: la parametrización del patrón mismo como una función de orden superior \cite{SICP}. Los tres patrones fundamentales —\eqfn{map}, \eqfn{filter} y \eqfn{fold}— no son funciones de biblioteca arbitrarias sino la manifestación de estructuras algebraicas profundas sobre el tipo $\eqdom{List}$.

\subsubsection*{\texttt{map}: Transformación Uniforme}

\eqfn{map} aplica una función $\eqfn{f} : \eqdom{A} \to \eqdom{B}$ a cada elemento de una lista de tipo $\eqdom{List}(\eqdom{A})$, produciendo una lista de tipo $\eqdom{List}(\eqdom{B})$ con la misma longitud \cite{SICP}. Es un \textit{morfismo de funtores}: transforma los contenidos preservando la estructura de la lista.

\begin{definicion}{\texttt{map} como morfismo estructural}
La definición recursiva de $\eqfn{map}$ sigue directamente el template de §\ref{subsec:4_3_3_recursion_estructural}. La implementación nativa extiende esta definición a múltiples listas: $\eqop{(}\eqfn{map}\ \equ{f}\ \equ{L_1} \cdots \equ{L_n}\eqop{)}$ aplica $\equ{f}$ a los elementos correspondientes de todas las listas en paralelo \cite{SICP}:
\begin{lstlisting}[language=Lisp, caption={Definición recursiva de map}]
(define (my-map f L)
  (if (null? L)
      '()
      (cons (f (car L))
            (my-map f (cdr L)))))
\end{lstlisting}
\end{definicion}

\begin{ejemplo}{Aplicaciones de \texttt{map}}
\begin{lstlisting}[language=Lisp, caption={Uso de map con distintas funciones}]
(map abs '(1 -2 3 -4 5 -6))
; -> (1 2 3 4 5 6)

(map (lambda (x) (* x x)) '(1 -3 -5 7))
; -> (1 9 25 49)

(map cons '(a b c) '(1 2 3))
; -> ((a . 1) (b . 2) (c . 3))

(map + '(1 2) '(3 4) '(5 6))
; -> (9 12)
\end{lstlisting}
\end{ejemplo}

\subsubsection*{\texttt{filter}: Selección Predicativa}

\eqfn{filter} retiene únicamente los elementos de una lista que satisfacen un predicado $\equ{p} : \eqdom{A} \to \eqdom{Bool}$, devolviendo una sublista de los elementos que pasan el criterio \cite{HtDP}. Nótese que \eqfn{setdiff} del §\ref{subsec:4_3_4_operaciones} es esencialmente un \eqfn{filter} con el predicado de no-pertenencia: $\equ{p}(\equ{x}) = \lnot\eqfn{member}(\equ{x},\ \equ{lis2})$.

\begin{definicion}{\texttt{filter} como selección estructural}
\begin{lstlisting}[language=Lisp, caption={Definición recursiva de filter}]
(define (filter p L)
  (cond
    [(null? L) '()]
    [(p (car L)) (cons (car L) (filter p (cdr L)))]
    [else        (filter p (cdr L))]))
\end{lstlisting}
\end{definicion}

\begin{ejemplo}{Filtrado de listas}
\begin{lstlisting}[language=Lisp, caption={Uso de filter para selección predicativa}]
(filter even? '(1 2 3 4 5 6))
; -> (2 4 6)

(filter (lambda (x) (> x 3)) '(1 5 2 7 3 8))
; -> (5 7 8)
\end{lstlisting}
\end{ejemplo}

\subsubsection*{\texttt{fold}: El Catamorfismo sobre Listas}

\eqfn{fold} (también denominado \eqfn{reduce} o \eqfn{accumulate}) es el patrón más general: recorre una lista aplicando una función binaria $\equ{f}$ a cada elemento y un valor acumulado, reduciendo la lista a un único valor \cite{SICP, HtDP}. Mientras \eqfn{map} y \eqfn{filter} producen listas, \eqfn{fold} puede producir cualquier tipo — es el mecanismo de \textit{eliminación} del tipo $\eqdom{List}$.

\begin{definicion}{\texttt{fold-right} y \texttt{fold-left}}
Existen dos variantes según el orden de asociación \cite{SICP}:
\begin{lstlisting}[language=Lisp, caption={Definición de fold-right y fold-left}]
;; fold-right: f(e1, f(e2, f(e3, init)))
(define (fold-right f init L)
  (if (null? L)
      init
      (f (car L) (fold-right f init (cdr L)))))

;; fold-left: f(f(f(init, e1), e2), e3)  [tail-recursive]
(define (fold-left f acc L)
  (if (null? L)
      acc
      (fold-left f (f acc (car L)) (cdr L))))
\end{lstlisting}
Nótese que \eqfn{fold-left} es naturalmente recursivo de cola, mientras que \eqfn{fold-right} no lo es. Para funciones $\equ{f}$ asociativas y conmutativas el resultado es idéntico; en general difieren tanto en el valor como en el costo espacial \cite{SICP}.
\end{definicion}

\subsection*{\texttt{map} y \texttt{filter} como casos especiales de \texttt{fold}}

El resultado algebraico central de esta sección es que \eqfn{map} y \eqfn{filter} son instancias particulares de \eqfn{fold-right}, lo que establece una jerarquía de expresividad: todo lo que puede computarse con \eqfn{map} o \eqfn{filter} puede expresarse como un \eqfn{fold} \cite{SICP}.

\begin{equation}
    \eqfn{map}\ \equ{f}\ \equ{L} \;\eqop{\equiv}\; \eqfn{fold\text{-}right}\ \eqop{(}\eqfn{\lambda}\ \eqop{(}\equ{x}\ \equ{acc}\eqop{)}\ \eqop{(}\eqfn{cons}\ \eqop{(}\equ{f}\ \equ{x}\eqop{)}\ \equ{acc}\eqop{)}\eqop{)}\ \eqk{'()}\ \equ{L}
\end{equation}

\begin{equation}
    \eqfn{filter}\ \equ{p}\ \equ{L} \;\eqop{\equiv}\; \eqfn{fold\text{-}right}\ \eqop{(}\eqfn{\lambda}\ \eqop{(}\equ{x}\ \equ{acc}\eqop{)}\ \eqop{(}\eqk{if}\ \eqop{(}\equ{p}\ \equ{x}\eqop{)}\ \eqop{(}\eqfn{cons}\ \equ{x}\ \equ{acc}\eqop{)}\ \equ{acc}\eqop{)}\eqop{)}\ \eqk{'()}\ \equ{L}
\end{equation}

\begin{ejemplo}{Verificación constructiva de las equivalencias}
\begin{lstlisting}[language=Lisp, caption={map y filter expresados mediante fold-right}]
(define (map-via-fold f L)
  (fold-right (lambda (x acc) (cons (f x) acc)) '() L))

(define (filter-via-fold p L)
  (fold-right (lambda (x acc)
                (if (p x) (cons x acc) acc))
              '() L))

(map-via-fold (lambda (x) (* x x)) '(1 2 3 4))
; -> (1 4 9 16)

(filter-via-fold even? '(1 2 3 4 5 6))
; -> (2 4 6)
\end{lstlisting}
\end{ejemplo}

\begin{ejemplo}{Conjunto potencia (\textit{power set}) mediante composición}
La generación del conjunto potencia $\mathcal{P}(S)$ ilustra la composición de \eqfn{map} con recursión estructural. La identidad recursiva $\mathcal{P}(S) = \mathcal{P}(\eqfn{cdr}(S)) \cup \{\{s\} \cup T \mid T \in \mathcal{P}(\eqfn{cdr}(S)),\; s = \eqfn{car}(S)\}$ se traduce directamente \cite{SICP}:
\begin{lstlisting}[language=Lisp, caption={Generación del conjunto potencia}]
(define (power-set set)
  (if (null? set)
      '(())
      (let ((power-set-of-rest (power-set (cdr set))))
        (append power-set-of-rest
                (map (lambda (subset)
                       (cons (car set) subset))
                     power-set-of-rest)))))

; (power-set '(1 2 3))
; -> (() (3) (2) (2 3) (1) (1 3) (1 2) (1 2 3))
\end{lstlisting}
El resultado tiene $2^{|S|}$ elementos, consistente con la cardinalidad del conjunto potencia. \eqfn{map} aquí no itera sobre los elementos de \equ{set} sino sobre los subconjuntos ya computados, actuando como un operador de extensión que antepone \eqfn{car}(\equ{set}) a cada subconjunto del resto \cite{SICP}.
\end{ejemplo}

\begin{fisicaconexion}{Operadores de transformación, proyección y contracción tensorial}
Los tres patrones de iteración funcional tienen análogos directos en álgebra lineal y física matemática. \eqfn{map} corresponde a la acción de un operador lineal $\hat{A}$ sobre los vectores de una base: $\hat{A}|\phi_i\rangle \to |\psi_i\rangle$, transformando el espacio vectorial preservando su dimensión. \eqfn{filter} es análogo a un proyector $\hat{P}$ sobre un subespacio: selecciona los vectores base que pertenecen al subespacio del predicado, reduciendo efectivamente la dimensión. \eqfn{fold} generaliza ambos: es una contracción tensorial que reduce un objeto de rango $n$ a uno de rango menor mediante la aplicación iterada de una forma bilineal — el mecanismo de la traza $\operatorname{Tr}(T^{ij}\,g_{jk})$ que reduce un tensor de segundo orden a uno de primero. La jerarquía \eqfn{map}, \eqfn{filter} $\subset$ \eqfn{fold} refleja la jerarquía de operaciones tensoriales ordenadas según su grado de contracción.
\end{fisicaconexion}


\section{Semántica y Terminación}
\label{sec:4_5_semantica}

Las secciones anteriores han construido el vocabulario operacional del paradigma funcional: cómo se definen funciones, cómo se procesan listas, cómo se abstraen patrones. Esta sección examina las garantías formales que subyacen a ese vocabulario — bajo qué condiciones un programa termina, qué forma adopta el proceso computacional que genera, y qué propiedad semántica asegura que los errores de tipo sean detectables y bien definidos \cite{SICP, PFPL}.

\subsection{Procesos Recursivos vs. Procesos Iterativos}
\label{subsec:4_5_1_recursion_vs_iteracion}

Una de las distinciones más importantes —y frecuentemente malentendidas— en la programación funcional es la diferencia entre un \textbf{proceso recursivo} y un \textbf{proceso iterativo}. Esta distinción no reside en la forma sintáctica del código, que en ambos casos utiliza \eqop{define} con auto-referencia, sino en la \textit{forma del proceso} que esa definición genera en ejecución \cite{SICP}.

Un \textbf{proceso recursivo} acumula \textit{trabajo diferido}: en cada llamada el marco de activación actual no puede completarse hasta que retorne la llamada interna, por lo que debe permanecer en la pila. La pila crece linealmente con la profundidad de la recursión y el proceso exhibe una fase de expansión seguida de una fase de contracción \cite{SICP}. Un \textbf{proceso iterativo} mantiene el estado computacional \textit{completamente} en los parámetros del acumulador en cada instante. No existe trabajo diferido: cuando se realiza la llamada recursiva, el marco actual ya ha completado todo su procesamiento y puede descartarse, manteniendo la pila de tamaño constante \cite{SICP, HtDP}.

\begin{ejemplo}{Factorial: proceso recursivo vs. proceso iterativo}
La versión directa genera un proceso recursivo con trabajo diferido:
\begin{lstlisting}[language=Lisp, caption={Factorial: proceso recursivo (expansión y contracción)}]
(define (factorial n)
  (if (= n 0)
      1
      (* n (factorial (- n 1)))))

; Expansion del proceso para n = 4:
; (factorial 4)
; (* 4 (factorial 3))
; (* 4 (* 3 (factorial 2)))
; (* 4 (* 3 (* 2 (factorial 1))))
; (* 4 (* 3 (* 2 (* 1 (factorial 0)))))
; (* 4 (* 3 (* 2 (* 1 1))))  <- contraccion
; 24
\end{lstlisting}

La versión con acumulador genera un proceso iterativo cuyo estado completo reside en \equ{factpartial}:
\begin{lstlisting}[language=Lisp, caption={Factorial: proceso iterativo (tail-recursive)}]
(define (facthelper n factpartial)
  (if (= n 0)
      factpartial
      (facthelper (- n 1) (* n factpartial))))

(define (factorial n) (facthelper n 1))

; Proceso para n = 4 (pila O(1)):
; (facthelper 4 1)
; (facthelper 3 4)
; (facthelper 2 12)
; (facthelper 1 24)
; (facthelper 0 24)  -> 24
\end{lstlisting}
\end{ejemplo}

La figura~\ref{fig:tail_vs_recursive} contrasta la forma de ambos procesos en términos del uso de la pila de llamadas.

\begin{figure}[htbp]
    \centering
    \begin{tikzpicture}[
        frame/.style={draw, rectangle, minimum width=4.0cm, minimum height=0.58cm,
                      fill=CtpBase, font=\footnotesize\ttfamily, align=center},
        pend/.style={draw, rectangle, minimum width=4.0cm, minimum height=0.58cm,
                     fill=CtpMantle, font=\footnotesize\ttfamily, align=center,
                     text=CtpSubtext0},
        >=latex
    ]
    %--- Proceso Recursivo (izquierda) ---
    \node[frame] (r0) at (0,0)    {(factorial 4)};
    \node[pend]  (r1) at (0,0.66) {(* 4 \textvisiblespace)\ \ \ \ \ \ \ \ \ };
    \node[pend]  (r2) at (0,1.32) {(* 3 \textvisiblespace)\ \ \ \ \ \ \ \ \ };
    \node[pend]  (r3) at (0,1.98) {(* 2 \textvisiblespace)\ \ \ \ \ \ \ \ \ };
    \node[pend]  (r4) at (0,2.64) {(* 1 \textvisiblespace)\ \ \ \ \ \ \ \ \ };
    \node[frame] (r5) at (0,3.30) {(factorial 0)};

    \draw[decorate, decoration={brace, amplitude=5pt, mirror},
          xshift=-2.3cm]
        (0,0.02) -- (0,3.58) node[midway, left=6pt, font=\scriptsize] {$O(n)$};

    \node[below=0.3cm of r0, font=\small\bfseries] {Proceso Recursivo};

    %--- Proceso Iterativo (derecha) ---
    \node[frame] (i0) at (7,0.00) {(facthelper 4\ \ 1)};
    \node[frame] (i1) at (7,0.80) {(facthelper 3\ \ 4)};
    \node[frame] (i2) at (7,1.60) {(facthelper 2\ 12)};
    \node[frame] (i3) at (7,2.40) {(facthelper 1\ 24)};
    \node[frame] (i4) at (7,3.20) {(facthelper 0\ 24)};

    \foreach \a/\b in {i0/i1, i1/i2, i2/i3, i3/i4}{
        \draw[->, dashed] (\a.north) -- (\b.south);
    }

    \draw[decorate, decoration={brace, amplitude=5pt},
          xshift=9.3cm]
        (0,0.02) -- (0,0.60) node[midway, right=6pt, font=\scriptsize] {$O(1)$};

    \node[below=0.3cm of i0, font=\small\bfseries] {Proceso Iterativo (TCO)};
    \end{tikzpicture}
    \caption{Proceso recursivo (izquierda): la pila crece $O(n)$ acumulando trabajo diferido (marcos en gris). Proceso iterativo con TCO (derecha): cada marco se descarta antes de la siguiente llamada, manteniendo pila $O(1)$ \cite{SICP}.}
    \label{fig:tail_vs_recursive}
\end{figure}

La razón por la que el proceso iterativo puede mantener una pila de tamaño constante es la \textbf{optimización de llamada de cola} (\textit{Tail-Call Optimization}, TCO). Una llamada es una \textit{tail call} si es la \textit{última operación} del procedimiento invocante — su resultado es directamente el resultado de la función sin ninguna computación pendiente. Cuando el intérprete detecta esta condición, descarta el marco de activación actual \textit{antes} de realizar la llamada, reutilizando el espacio de pila \cite{SICP, PFPL}.

\begin{observacion}{TCO como garantía semántica en Scheme}
El estándar R$^5$RS y sus sucesores R$^6$RS y R$^7$RS establecen que toda implementación conforme de \eqk{Scheme} \textit{debe} implementar TCO \cite{SICP}. Esto distingue a \eqk{Scheme} de lenguajes como \eqk{Python} o \eqk{Java}, donde la optimización de llamadas de cola no está garantizada y la recursión profunda puede producir un \textit{stack overflow}. En \eqk{Scheme}, un proceso iterativo expresado mediante recursión de cola tiene exactamente el mismo costo espacial que un bucle imperativo: $O(1)$ en la pila de llamadas.
\end{observacion}

\begin{fisicaconexion}{Transformación canónica y absorción de grados de libertad}
La conversión de recursión ordinaria a recursión de cola es análoga a una transformación canónica en mecánica hamiltoniana. En un sistema con $n$ grados de libertad, una transformación canónica bien elegida puede hacer cíclicos ciertos grados de libertad — sus momentos conjugados se convierten en constantes de movimiento que se «absorben» en los parámetros del sistema sin pérdida de información. El acumulador \equ{factpartial} desempeña exactamente ese rol: absorbe el trabajo pendiente como un momento conjugado que lleva toda la historia del cómputo, eliminando la necesidad de mantener la cadena de marcos activos en la pila.
\end{fisicaconexion}

\subsection{Mergesort: Recursión Estructural No Trivial}
\label{subsec:4_5_2_mergesort}

El algoritmo de ordenamiento por mezcla (\textit{merge sort}) es el caso de estudio paradigmático para demostrar que la recursión estructural escala con elegancia a algoritmos de complejidad asintótica óptima. La versión funcional sobre listas enlazadas presenta diferencias arquitectónicas respecto a la implementación imperativa clásica que merecen análisis explícito \cite{SICP}.

La estrategia divide la lista \equ{L} en dos sublistas aproximadamente iguales, las ordena recursivamente y las fusiona. Para listas enlazadas —que carecen de acceso aleatorio— la división se realiza por \textbf{separación intercalada}: \eqfn{oddsplit} extrae los elementos en posiciones impares y \eqfn{evensplit} los de posiciones pares. Ambas operaciones son $O(n)$, idéntico costo asintótico al de encontrar el punto medio en un arreglo \cite{SICP}.

\begin{definicion}{Implementación funcional de Mergesort}
\begin{lstlisting}[language=Lisp, caption={Mergesort funcional sobre listas enlazadas}]
;; Extrae elementos en posiciones impares (1, 3, 5, ...)
(define (oddsplit L)
  (cond
    [(null? L)        '()]
    [(null? (cdr L))  (list (car L))]
    [else (cons (car L) (oddsplit (cdr (cdr L))))]))

;; Extrae elementos en posiciones pares (2, 4, 6, ...)
(define (evensplit L)
  (cond
    [(null? L)        '()]
    [(null? (cdr L))  '()]
    [else (cons (car (cdr L)) (evensplit (cdr (cdr L))))]))

;; Fusion de dos listas ordenadas en O(n)
(define (mergelists L M)
  (cond
    [(null? L) M]
    [(null? M) L]
    [(< (car L) (car M))
     (cons (car L) (mergelists (cdr L) M))]
    [else
     (cons (car M) (mergelists L (cdr M)))]))

;; Mergesort principal
(define (mergesort L)
  (cond
    [(null? L)        L]
    [(null? (cdr L))  L]
    [else (mergelists (mergesort (oddsplit L))
                      (mergesort (evensplit L)))]))
\end{lstlisting}
\end{definicion}

\subsubsection*{Análisis de complejidad}

Sea $T(n)$ el número de operaciones elementales para ordenar una lista de $n$ elementos. \eqfn{oddsplit} y \eqfn{evensplit} producen sublistas de tamaño $\lceil n/2 \rceil$ y $\lfloor n/2 \rfloor$ respectivamente, cada una en tiempo $O(n)$. \eqfn{mergelists} fusiona dos listas ordenadas en tiempo $O(n)$. La recurrencia resultante es:

\begin{equation}
    T(n) = 2\,T\!\left(\frac{n}{2}\right) + O(n), \qquad T(1) = O(1)
\end{equation}

Por el \textbf{Teorema Maestro} con $a = 2$, $b = 2$, $f(n) = \Theta(n) = \Theta(n^{\log_2 2})$ — caso segundo — se obtiene $T(n) = O(n \log n)$, la cota inferior para comparaciones en el modelo algebraico de decisión \cite{PFPL}.

\begin{observacion}{El costo de la inmutabilidad en mergesort}
A diferencia del mergesort imperativo que opera sobre el mismo arreglo con un buffer auxiliar de tamaño $O(n)$, la versión funcional construye estructuras nuevas en cada llamada a \eqfn{cons}. El costo en memoria es $O(n \log n)$: en cada uno de los $O(\log n)$ niveles de recursión se construyen $O(n)$ cons-cells nuevas. La contrapartida es que las listas de entrada permanecen intactas — la inmutabilidad elimina el aliasing y permite que el mismo dato sea procesado concurrentemente por múltiples computaciones sin interferencia, una propiedad esencial en entornos paralelos \cite{SICP}.
\end{observacion}

\subsection{Modelo de Memoria y Recolección de Basura}
\label{subsec:4_5_3_gc}

La inmutabilidad del §\ref{subsec:4_1_2_transparencia} tiene una consecuencia inevitable: cada operación que «modifica» una estructura de datos produce en realidad una \textit{nueva} estructura sin alterar la original. En un programa activo, esto genera un flujo constante de objetos efímeros — creados, utilizados brevemente y abandonados. La \textbf{recolección de basura} (\textit{garbage collection}, GC) es el mecanismo que recupera automáticamente ese espacio sin intervención del programador \cite{SICP, HtDP}.

El criterio que determina si un objeto puede recolectarse es la \textbf{alcanzabilidad} (\textit{reachability}): un objeto es \textit{basura} si y solo si no existe ninguna cadena de referencias que lo conecte con las \textit{raíces} del entorno — la tabla de símbolos global y la pila de llamadas activa \cite{SICP}.

\begin{ejemplo}{Recolección por pérdida de alcanzabilidad}
Considérese la siguiente secuencia de definiciones \cite{SICP}:
\begin{lstlisting}[language=Lisp, caption={Pérdida de alcanzabilidad y recolección de basura}]
(define x '(1 2 3))  ; x apunta a cons-cell con car=1
(define y (cdr x))   ; y apunta al cons-cell con car=2
; ...
(define x '(4 5))    ; x ahora apunta a nueva lista
; El cons-cell con car=1 ya no es alcanzable desde
; ninguna raiz: es candidato a recoleccion.
; y aun referencia a (2 3), que sigue siendo alcanzable.
\end{lstlisting}
\end{ejemplo}

La figura~\ref{fig:gc_memory} ilustra el estado de la memoria tras la redefinición de \eqk{x}.

\begin{figure}[htbp]
    \centering
    \begin{tikzpicture}[
        cell/.style={rectangle split, rectangle split parts=2,
                     rectangle split horizontal, draw,
                     minimum height=0.85cm, fill=CtpBase, font=\footnotesize},
        gcell/.style={rectangle split, rectangle split parts=2,
                      rectangle split horizontal, draw, dashed,
                      minimum height=0.85cm, fill=CtpMantle, font=\footnotesize,
                      text=CtpOverlay1},
        symb/.style={draw, rectangle, fill=CtpSurface0, font=\small\bfseries,
                     minimum width=1.1cm, minimum height=0.55cm},
        >=latex, node distance=1.8cm
    ]

    % Tabla de simbolos
    \node[font=\small\bfseries, text=CtpSubtext0] at (-0.5, 4.0) {Tabla de símbolos};
    \node[symb] (sx) at (-0.5, 3.2) {\eqk{x}};
    \node[symb] (sy) at (-0.5, 1.7) {\eqk{y}};

    % Lista nueva de x: (4 5 nil)
    \node[cell] (c4) at (2.8, 3.2) {\eqk{4} \nodepart{two} };
    \node[cell] (c5) at (4.9, 3.2) {\eqk{5} \nodepart{two} };
    \node[draw, circle, font=\scriptsize, minimum size=0.5cm] (nil1) at (6.6, 3.2) {$\emptyset$};
    \draw[->] (sx.east) -- (c4.west);
    \draw[->] (c4.two east) -- (c5.west);
    \draw[->] (c5.two east) -- (nil1);

    % Lista de y: (2 3 nil) — alcanzable
    \node[cell] (c2) at (2.8, 1.7) {\eqk{2} \nodepart{two} };
    \node[cell] (c3) at (4.9, 1.7) {\eqk{3} \nodepart{two} };
    \node[draw, circle, font=\scriptsize, minimum size=0.5cm] (nil2) at (6.6, 1.7) {$\emptyset$};
    \draw[->] (sy.east) -- (c2.west);
    \draw[->] (c2.two east) -- (c3.west);
    \draw[->] (c3.two east) -- (nil2);

    % Cons-cell basura: car=1 (inalcanzable)
    \node[gcell] (c1) at (2.8, 0.2) {\eqk{1} \nodepart{two} };
    \draw[->, dashed] (c1.two east) to[out=0, in=270] (c2.south);
    \node[font=\scriptsize, text=CtpOverlay1, below=0.15cm of c1]
          {inalcanzable $\to$ GC};

    \end{tikzpicture}
    \caption{Estado de la memoria tras \texttt{(define x '(4 5))}. El cons-cell con \eqk{car=1} (líneas discontinuas) ya no es alcanzable desde ninguna raíz y es candidato a recolección. El cons-cell \eqk{car=2} sigue siendo alcanzable mediante \eqk{y} \cite{SICP}.}
    \label{fig:gc_memory}
\end{figure}

\begin{definicion}{Recolectores de Racket: generacional y conservador}
\eqk{Racket} implementa dos estrategias de recolección \cite{SICP, HtDP}:
\begin{itemize}
    \item \textbf{Generacional (3m)}: separa los objetos en generaciones según su tiempo de vida estimado. Los objetos jóvenes se recolectan frecuentemente en un \textit{nursery} de tamaño fijo; los que sobreviven se promueven a generaciones mayores, recolectadas con menor frecuencia. La \textit{hipótesis generacional} — que la mayoría de objetos muere joven — es especialmente válida en programación funcional, donde la inmutabilidad genera grandes cantidades de objetos intermedios de corta duración. Es la estrategia predeterminada en \eqk{Racket}.
    \item \textbf{Conservador (CGC)}: trata cualquier palabra de memoria que pueda interpretarse como un puntero como efectivamente un puntero, sin requerir metadatos de tipo precisos. Facilita la integración con código \eqk{C}, pero introduce imprecisión — puede retener como vivos objetos cuya representación numérica coincide con la de un puntero válido, aunque no sean alcanzables.
\end{itemize}
\end{definicion}

\begin{fisicaconexion}{Alcanzabilidad y decaimiento de estados excitados}
La alcanzabilidad en GC es análoga al criterio de estabilidad en sistemas dinámicos: un objeto es «estable» (vivo) si existe al menos una trayectoria desde una condición de frontera (las raíces) hasta él. Un objeto inalcanzable es como un estado excitado que ha perdido toda conexión con el estado base — decae espontáneamente y su energía (memoria) se libera al entorno. El recolector generacional explota el equivalente de la vida media: al igual que los isótopos se clasifican por su vida media para diseñar estrategias de gestión radiológica, los objetos se clasifican por su tiempo de vida esperado para minimizar el costo amortizado de la recolección.
\end{fisicaconexion}

\subsection{Semántica Estática y Seguridad de Tipos}
\label{subsec:4_5_4_type_safety}

La discusión hasta este punto ha operado en el nivel de la \textbf{semántica dinámica}: las reglas que describen cómo un programa \textit{se ejecuta}. La \textbf{semántica estática} complementa este cuadro describiendo propiedades que pueden establecerse \textit{antes} de la ejecución, mediante el análisis del texto del programa \cite{PFPL, TAPL}. El objeto central de la semántica estática es el \textbf{sistema de tipos}.

\begin{definicion}{Seguridad de Tipos (\textit{Type Safety})}
Un lenguaje es \textbf{seguro en tipos} si su sistema de tipos satisface dos propiedades complementarias, formuladas sobre la relación entre la semántica estática (tipado, $\equ{e} : \eqdom{\tau}$) y la semántica dinámica (reducción, $\equ{e} \to \equ{e'}$) \cite{PFPL}:
\begin{enumerate}
    \item \textbf{Preservación} (\textit{Preservation}): si $\equ{e} : \eqdom{\tau}$ y $\equ{e} \to \equ{e'}$, entonces $\equ{e'} : \eqdom{\tau}$. La evaluación no altera el tipo de una expresión.
    \item \textbf{Progreso} (\textit{Progress}): si $\equ{e} : \eqdom{\tau}$ y $\equ{e}$ es una expresión cerrada, entonces $\equ{e}$ es un valor o existe $\equ{e'}$ tal que $\equ{e} \to \equ{e'}$. Un programa bien tipado nunca queda atascado en un estado sin transición definida.
\end{enumerate}
Juntas, estas propiedades garantizan que la evaluación de un programa bien tipado o produce un valor del tipo correcto o diverge — pero nunca falla con comportamiento indefinido \cite{PFPL, TAPL}.
\end{definicion}

La distinción entre \textbf{tipado estático} y \textbf{tipado dinámico} no es sobre \textit{si} se verifican los tipos, sino sobre \textit{cuándo} \cite{PFPL, TAPL}. En lenguajes con tipado estático como \eqk{Haskell} o \eqk{ML}, el verificador analiza el texto del programa antes de la ejecución, rechazando en tiempo de compilación cualquier expresión que violaría type safety. En \eqk{Scheme}, la verificación ocurre en tiempo de ejecución: los tipos residen en los valores, no en las variables, y el intérprete comprueba la compatibilidad inmediatamente antes de cada operación \cite{SICP, HtDP}.

\begin{observacion}{Scheme es type-safe; C no lo es}
Es crucial no confundir tipado dinámico con ausencia de seguridad de tipos. \eqk{Scheme} es \textbf{type-safe}: la expresión \eqop{(}\eqfn{+}\ \eqk{1}\ \eqk{"hola"}\eqop{)} produce un error de tipo bien definido en tiempo de ejecución — nunca comportamiento indefinido \cite{SICP}. \eqk{C}, en contraste, no es type-safe: el \textit{cast} explícito permite reinterpretar arbitrariamente la representación binaria de cualquier valor, y el acceso fuera de los límites de un arreglo produce comportamiento indefinido sin diagnóstico garantizado. El precio del tipado dinámico en \eqk{Scheme} es que los errores se detectan en ejecución en lugar de compilación; la ganancia es la flexibilidad de trabajar con estructuras heterogéneas sin declaraciones de tipo a priori \cite{HtDP}.
\end{observacion}

La tensión fundamental entre ambas estrategias puede formularse en términos de la \textit{indecidibilidad}: el problema de determinar si un programa arbitrario es type-safe es, en general, indecidible. Los sistemas de tipos estáticos resuelven esto siendo \textit{conservadores} — rechazan programas que \textit{podrían} fallar, pero también algunos que son correctos. Esta sobre-aproximación es el costo de la detección temprana. Los lenguajes de tipado dinámico son precisos pero tardíos: solo rechazan lo que efectivamente falla en ejecución \cite{PFPL, TAPL}.

\begin{fisicaconexion}{Leyes de conservación como invariantes de tipo}
La type safety es el análogo computacional de una ley de conservación en física. La \textbf{Preservación} afirma que la «carga de tipo» $\eqdom{\tau}$ se conserva bajo la dinámica de evaluación — exactamente como la energía o el momento angular se conservan bajo la evolución hamiltoniana. El \textbf{Progreso} garantiza que el sistema no alcanza singularidades ni estados absorbentes inesperados: toda configuración bien tipada tiene una transición definida, análogamente a como las soluciones de ecuaciones de campo bien planteadas no desarrollan singularidades en tiempo finito bajo condiciones de contorno apropiadas. Un programa con un error de tipo es entonces análogo a un sistema con una ley de conservación violada — el comportamiento posterior es impredecible, potencialmente catastrófico, y sin garantía de diagnóstico.
\end{fisicaconexion}

\section*{Resumen del Capítulo}
\addcontentsline{toc}{section}{Resumen del Capítulo}

La programación funcional constituye un cambio de paradigma que desplaza el centro de gravedad del cómputo desde la manipulación de estado hacia la evaluación de expresiones matemáticas. El hilo conductor del capítulo puede resumirse en cinco tesis que se refuerzan mutuamente.

\textbf{La transparencia referencial es la propiedad fundacional.} Al garantizar que $\eqfn{f}(\equ{a})$ produce siempre el mismo valor para el mismo argumento $\equ{a}$, se hace posible razonar sobre el programa mediante sustitución de iguales por iguales — el mecanismo de la $\beta$-reducción del $\lambda$-calculus. Esta propiedad no es una restricción arbitraria sino la condición que habilita la verificación formal, la optimización segura y el paralelismo sin interferencias \cite{SICP, PFPL}.

\textbf{La estructura de los datos determina la estructura del código.} La definición inductiva de $\eqdom{List}$ — base $\eqk{'()}$ más caso constructor \eqfn{cons} — se refleja directamente en el patrón de recursión estructural: caso base \eqfn{null?} más caso inductivo sobre \eqfn{car}/\eqfn{cdr}. Esta correspondencia, sistematizada por HtDP como receta de diseño, garantiza la cobertura exhaustiva de los casos y proporciona una ruta directa hacia la demostración de terminación \cite{HtDP, PFPL}.

\textbf{Las funciones como ciudadanos de primera clase generan una jerarquía de abstracción.} La capacidad de pasar funciones como argumentos y retornarlas como resultados permite capturar patrones computacionales recurrentes — \eqfn{map}, \eqfn{filter}, \eqfn{fold} — que de otro modo requerirían duplicación de código. El resultado algebraico central es que \eqfn{map} y \eqfn{filter} son instancias particulares de \eqfn{fold-right}, estableciendo una jerarquía de expresividad sobre el tipo $\eqdom{List}$ \cite{SICP}.

\textbf{La forma del proceso no coincide necesariamente con la forma del código.} Un procedimiento con auto-referencia sintáctica puede generar un proceso recursivo con pila $O(n)$ o un proceso iterativo con pila $O(1)$, dependiendo de si existe trabajo diferido tras la llamada recursiva. La optimización de llamada de cola (TCO), garantizada como propiedad semántica por el estándar R$^5$RS, hace que la recursión de cola sea equivalente en costo espacial a un bucle imperativo, eliminando la dicotomía artificial entre recursión e iteración \cite{SICP, PFPL}.

\textbf{La type safety es la garantía formal de ausencia de comportamiento indefinido.} La preservación — los tipos se conservan bajo evaluación — y el progreso — un programa bien tipado nunca queda atascado — son las dos propiedades que, juntas, aseguran que un programa en \eqk{Scheme} o bien termina en un valor bien definido, o bien diverge, pero nunca falla silenciosamente. El tipado dinámico de \eqk{Scheme} satisface ambas propiedades verificando en tiempo de ejecución, lo que lo distingue de \eqk{C} — que no es type-safe — y lo sitúa en el mismo espacio de garantías formales que los lenguajes con tipado estático, con la diferencia de \textit{cuándo} se realiza la verificación \cite{PFPL, TAPL}.

\begin{observacion}{Lenguajes funcionales más allá de Scheme}
\eqk{Scheme} es el vehículo pedagógico de este capítulo por su minimalismo semántico, pero el paradigma funcional se extiende a un ecosistema amplio. \eqk{Haskell} incorpora tipado estático con inferencia de tipos (sistema Hindley-Milner), evaluación perezosa (\textit{lazy evaluation}) y un sistema de clases de tipos que generaliza la sobrecarga. \eqk{ML} y \eqk{OCaml} combinan tipado estático con efectos controlados. \eqk{Erlang} añade la concurrencia masiva mediante procesos ligeros sin estado compartido. En todos los casos, las ideas de transparencia referencial, funciones de orden superior y recursión estructural desarrolladas aquí constituyen el núcleo conceptual compartido \cite{SICP, HtDP}.
\end{observacion}

\printbibliography[segment=\therefsegment, heading=subbibliography]